\documentclass{article}
\usepackage{iclr2024_conference,times}
\input{math_commands.tex}
\usepackage{booktabs}
\usepackage{hyperref}
\usepackage{url}
\usepackage{graphicx}
\usepackage{multirow}
\usepackage{array}
\usepackage{colortbl}
\usepackage{enumitem}
\iclrfinalcopy

\title{Study Coach: Report 3. Analysis and Initial Modifications}
\author{
  Dipti Aswath$^*$ \hspace{2em} Buddhika Makehelwala$^*$ \hspace{2em} Jiashan Cui$^*$ \hspace{2em} Khubi Shah$^*$ \\
  \texttt{\{daswath, bmakehelwala, jiashanc, khubis\}@andrew.cmu.edu}
}
\date{}

\begin{document}
\maketitle

The aim in this report is to analyze the outputs of the approaches collected in Report 2,
and try out modifications to see how they affect performance on intrinsic metrics
and qualitatively. Our task is incongruence detection: given a scientific figure and
a student answer, a model must (1) classify the verdict (correct / partially correct
/ incorrect), (2) identify the error category (factual, conceptual, omission), and
(3) generate coaching feedback.

\section{Analysis (4 pages)}

\subsection{\points{2} Intrinsic Metrics}

\noindent\textit{\textbf{NOTE:} This subsection is a WIP and will be refined.}
\medskip

For a VLM to grade student answers about scientific figures, it needs three core
skills: recognizing incomplete answers, knowing when coaching is required, and
extracting specific values from visual data. We measure these independently of
end-to-end verdict accuracy using per-example outputs from the 8B evaluation set.

SPIQA+ contains 174 examples with three verdict classes: 66 correct (37.9\%),
61 partially correct (35.1\%), and 47 incorrect (27.0\%). Since coaching is only
needed for non-correct answers, we compute intrinsic metrics on the 108 examples
where ground truth is partially correct or incorrect.

\paragraph{Intrinsic Metric 1: Partially-Correct Recall (PC Recall).}

Can the model recognize that a student answer is \textit{incomplete} rather than
fully wrong or fully right? Study Coach's core coaching scenario involves answers
that are partially correct---the student understood something but missed something
else. An educator who cannot reliably identify this middle class cannot target
feedback appropriately.

We define PC Recall as the fraction of ground-truth Partially Correct examples
correctly classified:
\begin{equation}
\text{PC Recall} = \frac{\text{TP}_\text{PC}}{\text{TP}_\text{PC} + \text{FN}_\text{PC}}
\end{equation}

This is not overall verdict accuracy: a model can maintain moderate accuracy by
collapsing partially correct answers into the incorrect class, generating feedback
of the wrong kind. PC Recall isolates the hardest, most coaching-relevant class
and is computed on the 61 ground-truth partially correct examples only.

At 8B scale, PC Recall drops sharply as visual context is added: text-only
44.3\%~$\to$~multimodal 19.7\% (Table~\ref{tab:intrinsic}). The model correctly
identifies fewer than 1 in 5 partially correct answers in the multimodal
condition---the same condition that produces the best coaching feedback. This
sharpens the verdict--feedback dissociation from Report 2 into a precise, per-class
finding: visual input specifically degrades recognition of the middle class that
most needs coaching.

\paragraph{Intrinsic Metric 2: Feedback Suppression Rate (FSR).}

Does the model generate coaching when it is needed? When the model predicts
``Correct,'' it outputs N/A feedback by prompt design. Suppressing feedback on
a truly correct answer is appropriate; suppressing it on an incorrect or partially
correct answer silences the coaching pipeline when the student needs help. We define
FSR on the 108 non-correct examples only:
\begin{equation}
\text{FSR} = \frac{\text{\# Correct predictions on non-correct GT}}{\text{\# non-correct GT examples (}N{=}108\text{)}}
\end{equation}

FSR measures a prerequisite skill: recognizing that coaching is required at all.
A model with high FSR produces false-positive verdicts with no explanation,
directly harming the student. FSR is distinct from verdict accuracy---a model
could classify incorrect answers correctly while still suppressing feedback on
every partial-credit case.

FSR increases substantially with visual context: text-only 13.9\%~$\to$~vision-only
30.6\%~$\to$~multimodal 28.7\% (Table~\ref{tab:intrinsic}). Adding visual input
causes the model to incorrectly accept more student answers as correct,
particularly for omission errors where the answer is plausible but structurally
incomplete. In the text-only condition, the model inappropriately suppresses
feedback on roughly 1 in 7 non-correct examples; in visual conditions, this rises
to nearly 1 in 3.

\paragraph{Intrinsic Metric 3: Visual Grounding Score (VGS).}

Can the model extract specific values from figures and reference them in feedback?
All three error types benefit from visual grounding (Report 2, H3): factual errors
require verifying specific numbers (``The figure shows 84.6\%, not 89\%''),
conceptual errors benefit from citing values that reveal patterns (``X increased
from 10 to 50''), and omission errors---which show the \textit{largest} feedback
gains (+33pp)---require the figure to reveal what the student failed to mention.
We define VGS on all 108 non-correct examples, excluding suppressed feedbacks:
\begin{equation}
\text{VGS} = \frac{\text{\# feedbacks citing specific values from figure}}{\text{\# non-suppressed feedbacks (}N{=}108\text{)}}
\end{equation}

VGS measures whether visual information is actually translated into text during
reasoning, independently of whether the verdict is correct. It is conceptually
related to Qwen3-VL's square-root reweighting loss, which scales multimodal
sample loss by $\sqrt{N_{\text{text}}/N_{\text{multimodal}}}$ to preserve
visual grounding during pre-training without degrading text capabilities.
VGS captures whether this grounding carries through to inference on our task.

VGS increases substantially with visual input: text-only 35.5\%~$\to$~vision-only
52.0\%~$\to$~multimodal 49.4\% at 8B (Table~\ref{tab:intrinsic}). At 32B, the
pattern strengthens: C1 37.2\%~$\to$~C4 56.0\% ($\Delta{=}+$18.8pp). This confirms
that the model \textit{does} use visual information for explanation generation even
when it cannot use it for classification---providing intrinsic evidence for the
verdict--feedback dissociation. Omission errors show the largest VGS gains
(+19.2pp at 8B), consistent with Report 2's finding that the figure reveals what
the student failed to mention.

\medskip
\noindent
\begin{center}
\small
\begin{tabular}{lll}
\toprule
Intrinsic Metric & Skill Measured & Why It Is Needed \\
\midrule
PC Recall & Detecting incomplete answers (n=61 PC) & Prerequisite for targeted coaching \\
FSR       & Knowing when coaching is needed (n=108) & Prerequisite for feedback generation \\
VGS       & Extracting values from figures (n=108)  & Prerequisite for grounded correction \\
\bottomrule
\end{tabular}
\end{center}

\paragraph{Evaluation methodology: LLM-as-Judge for feedback quality.}

Because open-ended coaching feedback cannot be evaluated with n-gram metrics alone
(F1, ROUGE-L, and BLEU show near-zero variation across conditions at both scales;
see Report 2, Table~10), we use Claude Sonnet~4 as an automated judge. Each predicted
feedback is classified as \textbf{Match} (correctly identifies error and explains
why), \textbf{Partial} (right direction but imprecise or incomplete), or
\textbf{Unmatched} (does not align with ground truth). Full definitions are in Report 2.

We report \textbf{Match\%} (strict: exact matches only) and \textbf{Soft Match\%}
(Match + Partial). At 32B, C4 achieves 20\% strict match but 58\% soft
match---meaning 38\% of responses are approximately correct but not precise.
Human evaluation ($n{=}10$/condition, 8B) confirmed the pattern: human match
ranged from 40\% (C1) to 80\% (C4), validating soft match as the practical metric
for coaching. We use strict match (Match\%) in Table~\ref{tab:intrinsic} since it
separates conditions more cleanly.

\clearpage

\begin{table}[t]
\centering
\scriptsize
\setlength{\tabcolsep}{4pt}
\begin{tabular}{@{}llrrrrrr@{}}
\toprule
& & \multicolumn{3}{c}{Intrinsic Metrics} & \multicolumn{3}{c}{Extrinsic Metrics} \\
\cmidrule(lr){3-5}\cmidrule(lr){6-8}
Model & Condition
  & PC Rec $\uparrow$
  & FSR $\downarrow$
  & VGS $\uparrow$
  & Verd $\uparrow$
  & Match\% $\uparrow$
  & SoftM\% $\uparrow$ \\
\midrule
\multicolumn{7}{l}{\textit{Unimodal Baselines (Qwen3-VL-8B, N=50 extrinsic, N=108 intrinsic)}} \\
Qwen3-VL-8B & C1: text-only    & 44.3\% & 13.9\% & 35.5\% & 56.0\% & 2.0\% & 40.0\% \\
Qwen3-VL-8B & C2: caption-only & 29.5\% & 17.6\% & 38.2\% & 48.0\% & 2.0\% & 46.0\% \\
\midrule
\multicolumn{7}{l}{\textit{Simple Multimodal Baselines (Qwen3-VL-8B, N=50 extrinsic, N=108 intrinsic)}} \\
Qwen3-VL-8B & C3: vision-only  & 21.3\% & 30.6\% & 52.0\% & 50.0\% & 4.0\% & 48.0\% \\
Qwen3-VL-8B & C4: multimodal   & 19.7\% & 28.7\% & 49.4\% & 48.0\% & 6.0\% & 56.0\% \\
\midrule
\multicolumn{7}{l}{\textit{Competitive Approaches (Qwen3-VL-32B, N=50 extrinsic, N=108 intrinsic)}} \\
Qwen3-VL-32B & C1: text-only    & \textit{TBD} & \textit{TBD} & 37.2\% & 56.0\% & 2.0\% & 42.0\% \\
Qwen3-VL-32B & C2: caption-only & \textit{TBD} & \textit{TBD} & 45.3\% & 58.0\% & 8.0\% & 48.0\% \\
Qwen3-VL-32B & C3: vision-only  & \textit{TBD} & \textit{TBD} & 48.5\% & 56.0\% & 6.0\% & 48.0\% \\
Qwen3-VL-32B & C4: multimodal   & \textit{TBD} & \textit{TBD} & 56.0\% & 54.0\% & 20.0\% & 58.0\% \\
\midrule
\multicolumn{7}{l}{\textit{Modifications (8B and 32B, N=50 extrinsic, N=108 intrinsic)}} \\
Routing (8B)  & C1 verdict + C4 fdbk       & 44.3\% & 28.7\% & 49.4\% & 56.0\% & 6.0\% & 56.0\% \\
Routing (32B) & C1 verdict + C4 fdbk       & \textit{TBD} & \textit{TBD} & 56.0\% & 56.0\% & 20.0\% & 58.0\% \\
Qwen3.5-9B   & C1: text-only (Mod~2)  & \textit{TBD} & \textit{TBD} & \textit{TBD} & \textit{TBD} & \textit{TBD} & \textit{TBD} \\
Qwen3.5-9B   & C4: multimodal (Mod~2) & \textit{TBD} & \textit{TBD} & \textit{TBD} & \textit{TBD} & \textit{TBD} & \textit{TBD} \\
Qwen3.5-397B & C1: text-only (Mod~3)  & \textit{TBD} & \textit{TBD} & \textit{TBD} & \textit{TBD} & \textit{TBD} & \textit{TBD} \\
Qwen3.5-397B & C4: multimodal (Mod~3) & \textit{TBD} & \textit{TBD} & \textit{TBD} & \textit{TBD} & \textit{TBD} & \textit{TBD} \\
\bottomrule
\end{tabular}
\caption{\points{1} Intrinsic and extrinsic metrics across conditions and models.}
\label{tab:intrinsic}
\end{table}

\noindent\textbf{Reading Table~\ref{tab:intrinsic}:}

\textit{Intrinsic metrics} (computed from per-example outputs, N=108; 8B complete, 32B TBD):
\begin{itemize}[itemsep=1pt, topsep=2pt, leftmargin=*]
  \item \textbf{PC Recall} $\uparrow$ --- Partially-correct recall: fraction of ground-truth partial examples correctly classified
  \item \textbf{FSR} $\downarrow$ --- Feedback suppression rate: fraction of predictions with N/A feedback (lower is better)
  \item \textbf{VGS} $\uparrow$ --- Visual grounding score: fraction of non-suppressed feedbacks citing specific numerical values (all error types)
\end{itemize}

\textit{Extrinsic metrics} (from H1/H2 evaluation, N=50 per condition):
\begin{itemize}[itemsep=1pt, topsep=2pt, leftmargin=*]
  \item \textbf{Verdict} $\uparrow$ --- Three-class classification accuracy (correct / partially correct / incorrect)
  \item \textbf{Match\%} $\uparrow$ --- LLM-as-Judge (Claude Sonnet 4) strict match: feedback correctly identifies error and explains why
  \item \textbf{SoftM\%} $\uparrow$ --- Soft match (Match + Partial): feedback at least partially useful for coaching
\end{itemize}

\noindent\textit{Match categories (see Report 2 for full definitions):}
\begin{itemize}[itemsep=1pt, topsep=2pt, leftmargin=*]
  \item \textbf{Match} --- feedback aligns with ground truth: correctly identifies what is wrong and explains why
  \item \textbf{Partial} --- captures some right reasoning but misses key elements, or right direction but imprecise
  \item \textbf{Unmatched} --- feedback does not align with ground truth
\end{itemize}

\textit{Key observations (8B):}
\begin{itemize}[itemsep=1pt, topsep=2pt, leftmargin=*]
  \item \textbf{PC Recall drops with visual input:} C1 (44.3\%) $\to$ C4 (19.7\%)---visual information introduces noise into classification boundary
  \item \textbf{VGS rises with visual input:} C1 (35.5\%) $\to$ C4 (49.4\%)---model extracts more values from figures into feedback
  \item \textbf{Same visual input, opposite effects:} Hurts classification ($-$8pp verdict) but helps explanation ($+$4pp Match)
  \item \textbf{Routing recovers both:} C1 verdict (56\%, 44.3\% PC Recall) + C4 feedback (6\% Match, 49.4\% VGS)
\end{itemize}

\textit{Key observations (32B --- TBD):}
\begin{itemize}[itemsep=1pt, topsep=2pt, leftmargin=*]
  \item \textbf{PC Recall with visual input:} C1 (\textit{TBD}) $\to$ C4 (\textit{TBD})---does scaling reduce the drop?
  \item \textbf{VGS with visual input:} C1 (\textit{TBD}) $\to$ C4 (\textit{TBD})---does scaling improve grounding?
  \item \textbf{FSR comparison:} Does 32B suppress less feedback than 8B?
\end{itemize}

\textit{Scaling benefit from Report 2 (8B $\to$ 32B extrinsic metrics):}
\begin{itemize}[itemsep=1pt, topsep=2pt, leftmargin=*]
  \item \textbf{Verdict gap narrows:} $-$8pp (8B) $\to$ $-$2pp (32B)---visual input less harmful at scale
  \item \textbf{Feedback match improves:} 6\% (8B) $\to$ 20\% (32B) at C4---$3.3{\times}$ improvement
  \item \textbf{Figure type matters:} Schematics easiest to explain (46\%), tables hardest (12\%) despite being easiest to classify
  \item \textbf{Error type matters:} Omission errors 0\% at C1 $\to$ 33\% at C4---figure reveals what student omitted
\end{itemize}

\subsection{Modifications and their Effects}
\points{2}

\noindent\textit{\textbf{NOTE:} This subsection is a WIP and will be refined.}
\medskip

We implement three modifications motivated by the qualitative failure analysis in
Section~\ref{sec:qual} and the intrinsic metric patterns above. Modification~1
(Condition Routing) is a zero-cost inference-time fix computable from existing Report 2
outputs. Modifications~2 and 3 are architectural ablations that test competing
explanations for the verdict--feedback dissociation by swapping the model while
holding the prompt and evaluation pipeline constant.

\paragraph{Modification 1: Condition Routing (C1 verdict + C4 feedback).}

\textbf{Motivation.} The intrinsic metrics make the verdict--feedback dissociation
precise: PC Recall drops from 44.3\% to 19.7\% as visual context is added (C1 to
C4), while VGS rises from 35.5\% to 49.4\%. These are not correlated failures of
a single configuration---they measure different sub-skills that are genuinely in
tension at 8B scale. The model can use visual information for extraction (improving
VGS and feedback quality) but not for classification (degrading PC Recall). Routing
assigns each sub-task to the condition best suited to it: C1 for verdict and error
category, C4 for coaching feedback.

\textbf{Implementation.} No new inference is required. We pair per-example C1
predictions (verdict, error category) with per-example C4 predictions (feedback)
using matched (\texttt{paper\_id}, \texttt{question}) keys across the H3/H4
evaluation set. All 108 examples matched successfully (100\% coverage).

\textbf{Results.} Routing achieves the best combined profile across all conditions:
verdict 56.0\% and PC Recall 44.3\% (matching C1), VGS 49.4\% and Fdbk Match
6.0\% (matching C4), with FSR at 28.7\%. The routing system recovers C1's verdict
precision and C4's visual grounding simultaneously at zero additional inference
cost. The FSR limitation (28.7\%) arises because N/A feedback comes from C4 when
C4 predicts Correct; a production system would override this by using the C1
verdict to force feedback generation regardless of C4's verdict.

\paragraph{Modification 2: Architectural Ablation --- Qwen3.5-9B Early Fusion vs. Qwen3-VL-8B Bolt-On ViT.}

\textbf{Motivation.} The intrinsic metrics reveal a structural paradox: VGS rises
from 35.5\% (C1) to 49.4\% (C4), confirming that visual information is being
extracted and used for explanation generation---yet PC Recall simultaneously
collapses from 44.3\% to 19.7\%, and FSR roughly doubles. The model \emph{can}
use figure content; it just cannot apply it to the verdict classification decision.

One architectural explanation is the \emph{bolt-on} design of Qwen3-VL.
Per the Qwen3-VL technical report, the model uses a three-module pipeline: a
SigLIP-2 SO-400M vision encoder (trained separately from the LLM), a two-layer
MLP merger augmented by DeepStack residual connections (which inject features
from three intermediate ViT layers into the first three LLM layers), and a
Qwen3-8B/32B decoder backbone. The modular separation is confirmed by the
pre-training recipe: Stage~0 freezes both the vision encoder and the LLM
backbone while training \emph{only the merger}, establishing a clear
representational boundary between the two modalities. DeepStack improves
multi-level visual feature injection but does not eliminate this boundary;
the LLM backbone was pre-trained as a text model and receives visual tokens
as an external input stream rather than being co-trained with visual objectives
from the start. This late-injection design may explain why the model can use
injected visual features for explanation generation (improving VGS and feedback
quality) but cannot reliably route those features into the verdict decision
boundary (degrading PC Recall).

Qwen3.5-9B uses a fundamentally different design: \emph{early fusion}. Rather
than a three-module pipeline (separate SigLIP-2 encoder $\to$ MLP merger
$\to$ LLM), Qwen3.5-9B co-trains vision and language from the start of
pre-training within a unified Gated DeltaNet hybrid architecture---no separate
vision encoder, no MLP merger, no Stage~0 alignment phase. Visual and textual
representations occupy the same token space throughout all 32 layers from the
first training step. This is not a cosmetic difference: the model was never in
a regime where visual tokens arrived as an external stream bolted onto a
text-trained LLM backbone. If the three-module bolt-on design of Qwen3-VL
(SigLIP-2 encoder + MLP/DeepStack merger + frozen-then-thawed Qwen3 decoder)
causes the PC Recall collapse observed under visual input, Qwen3.5-9B should
show a smaller or reversed verdict gap (C4 accuracy $\geq$ C1 accuracy) at
comparable parameter count ($\sim$9B).

\textbf{Architecture comparison:}

\begin{center}
\small
\begin{tabular}{p{0.22\linewidth}p{0.32\linewidth}p{0.32\linewidth}}
\toprule
 & Qwen3-VL-8B (baseline) & Qwen3.5-9B (Mod~2) \\
\midrule
Vision encoder  & SigLIP-2 SO-400M (trained separately from LLM) & No separate encoder---early fusion \\
Cross-modal bridge & 2-layer MLP + DeepStack (3 ViT layers $\to$ 3 LLM layers) & Unified token space throughout \\
Pre-training    & Stage~0: merger only (ViT + LLM frozen); then staged full training & Co-trained from scratch \\
LLM backbone    & Qwen3-8B (standard transformer) & Gated DeltaNet + Gated Attention hybrid \\
Parameters      & 8B & 9B \\
Together.ai ID  & \texttt{Qwen/Qwen3-VL-8B-Instruct} & \texttt{Qwen/Qwen3.5-9B} \\
\bottomrule
\end{tabular}
\end{center}

\textbf{Exact implementation steps (what to run):}

\begin{enumerate}

\item \textbf{Confirm Together.ai access.}
Verify \texttt{Qwen/Qwen3.5-9B} is available under the existing Together.ai API
key used for Report 2. Model is live as of March 2, 2026 (\$0.10/1M input,
\$0.15/1M output tokens).

\item \textbf{Disable thinking mode.}
Qwen3.5-9B runs in thinking mode by default, which prepends a chain-of-thought
reasoning block before the final response. This will break the structured output
parser (which expects \texttt{Verdict = ...}, \texttt{Error Category = ...},
\texttt{Feedback = ...}). Disable it explicitly:
\begin{quote}
\texttt{extra\_body=\{"chat\_template\_kwargs": \{"enable\_thinking": False\}\}}
\end{quote}
Use sampling parameters: \texttt{temperature=0.0} (greedy, for reproducibility),
\texttt{top\_p=1.0}, \texttt{presence\_penalty=0.0}.

\item \textbf{Use identical prompts.}
Use the exact same prompt templates as Report 2 for all four conditions (C1--C4). No
prompt changes. The only change is the model ID string. This ensures any
performance difference is attributable to the model architecture, not prompt
variation.

\item \textbf{Run C1 and C4 on the H3/H4 evaluation set (N=108).}
These are the two conditions that define the verdict--feedback gap. C1 (text-only)
establishes the same-architecture text baseline. C4 (multimodal) is the
hypothesis-test condition: does early fusion eliminate the verdict drop?
Running all four conditions (C1--C4) is preferred if API budget allows, to produce
a complete comparison row in Table~\ref{tab:intrinsic}.

\item \textbf{Parse outputs with the same extraction logic.}
Qwen3.5-9B output formatting may differ slightly from Qwen3-VL (e.g., extra
whitespace, different capitalization of ``Partially Correct'' vs ``partially
correct''). Before computing metrics, verify the parser correctly extracts Verdict,
Error Category, and Feedback fields from 5--10 sample outputs. Normalize verdict
strings to lowercase before comparison.

\item \textbf{Run the same three evaluation scripts:}
  \begin{itemize}
  \item Verdict accuracy: compare predicted vs.\ ground-truth verdict labels
  \item Intrinsic metrics (PC Recall, FSR, VGS): same computation from raw
    per-example outputs
  \item Feedback quality (LLM-as-Judge): run \texttt{analyze\_feedback\_by\_error\_type.py}
    with \texttt{model=claude-sonnet-4-6}, \texttt{temperature=0}, N=108
  \end{itemize}

\item \textbf{Interpret the result against two hypotheses:}
  \begin{itemize}
  \item \textbf{If C4 verdict accuracy $\geq$ C1 for Qwen3.5-9B:} early fusion
    resolves the architectural mismatch. The Report 4 implication is to use Qwen3.5-9B
    as the backbone for fine-tuning on SPIQA+, not Qwen3-VL.
  \item \textbf{If C4 verdict accuracy $<$ C1 for Qwen3.5-9B:} the bolt-on ViT
    is not the cause. The failure is inherent to how visual input interacts with
    three-class verdict decisions at this parameter scale, and the Report 4 strategy
    shifts to fine-tuning Qwen3-VL-32B (or 72B+) with a Verdict-Classification
    Head.
  \end{itemize}

\end{enumerate}

\textbf{Results.} \textit{[To be filled after running N=108 on C1 and C4.
Update Table~\ref{tab:intrinsic} with PC Recall, FSR, VGS, Verdict, and Fdbk
Match rows for Qwen3.5-9B C1 and C4.]}

% FIGURE PLACEHOLDER (Mod 2): Grouped bar chart comparing Qwen3-VL-8B C1,
% Qwen3-VL-8B C4, Routing, Qwen3.5-9B C1, Qwen3.5-9B C4.
% Three clusters (PC Recall, FSR, Fdbk Match), five bars each.

\paragraph{Modification 3: Frontier-Scale Early Fusion --- Qwen3.5-397B-A17B (C1 and C4).}

\textbf{Motivation.} Modification~2 tests whether the bolt-on ViT architecture
causes the verdict--feedback dissociation by replacing it with early fusion at
matched scale ($\sim$9B). But a confound remains: perhaps the failure is not
architectural but simply a function of cross-modal reasoning capacity, which
scales with model size. Qwen3-VL-32B already narrowed the verdict gap from
$-$8pp to $-$2pp without resolving it. The question for Modification~3 is:
does early fusion at \emph{frontier scale} (17B active parameters, 397B total)
close the gap entirely?

This matters for Report 4 for a distinct reason from Modification~2. If Qwen3.5-9B
(Mod~2) resolves the gap, early fusion at 9B is sufficient and fine-tuning should
target that backbone. If Qwen3.5-9B does not resolve it but Qwen3.5-397B-A17B
does, the interaction of early fusion \emph{and} scale is required, and the Report 4
fine-tuning backbone must be at least the 397B MoE model (or the 27B dense
variant as a practical compromise). If neither resolves the gap, the failure is
inherent to the task at current scales, and the Report 4 strategy must rely on the
Verdict-Classification Head and routing rather than a single model.

The 397B-A17B is the same Qwen3.5 architecture as Modification~2 (early fusion,
Gated DeltaNet + MoE, unified token space), but scales to 397B total / 17B active
parameters via sparse mixture-of-experts with 512 total experts (10 routed + 1
shared per token).

\textbf{Architecture comparison across all three models:}

\begin{center}
\scriptsize
\begin{tabular}{p{0.18\linewidth}p{0.25\linewidth}p{0.22\linewidth}p{0.22\linewidth}}
\toprule
 & Qwen3-VL-8B (baseline) & Qwen3.5-9B (Mod~2) & Qwen3.5-397B-A17B (Mod~3) \\
\midrule
Vision encoder  & SigLIP-2 SO-400M (bolt-on, separate from LLM) & None---early fusion & None---early fusion \\
Cross-modal bridge & 2-layer MLP + DeepStack & Unified token space & Unified token space + MoE \\
LLM backbone    & Qwen3-8B (std transformer) & Gated DeltaNet hybrid & Gated DeltaNet + sparse MoE \\
Active params   & 8B                & 9B               & 17B \\
Total params    & 8B                & 9B               & 397B \\
Context         & 256K              & 262K             & 262K \\
Together.ai ID  & \texttt{Qwen3-VL-8B-Instruct} & \texttt{Qwen3.5-9B} & \texttt{Qwen3.5-397B-A17B} \\
\bottomrule
\end{tabular}
\end{center}

\textbf{Exact implementation steps (what to run):}

\begin{enumerate}

\item \textbf{Confirm Together.ai access and cost estimate.}
Model ID: \texttt{Qwen/Qwen3.5-397B-A17B}. Image input is supported natively.
At Together.ai rates, running N=108 examples across C1 and C4 (each example
$\approx$1,500 input + 300 output tokens; C4 adds $\approx$500 image tokens per
example) costs under \$5 total for both conditions combined.

\item \textbf{Disable thinking mode --- critical.}
Qwen3.5-397B-A17B runs in thinking mode by default, generating a long internal
reasoning block before the final response. This will break the structured output
parser. Disable it in every API call:
\begin{quote}
\texttt{extra\_body=\{"chat\_template\_kwargs": \{"enable\_thinking": False\}\}}
\end{quote}
Use \texttt{temperature=0.0}, \texttt{top\_p=1.0}, \texttt{presence\_penalty=0.0}
for reproducibility. If you prefer a small amount of temperature for diversity,
use \texttt{temperature=0.0} for the first run and compare to the 9B run, which
should also use \texttt{temperature=0.0}.

\item \textbf{Use identical prompts as Report 2 and Modification~2.}
No prompt changes from the Report 2 baseline. The only variable across all three
models (Qwen3-VL-8B, Qwen3.5-9B, Qwen3.5-397B-A17B) is the model ID string.
This is essential: comparing across models with different prompts would confound
the architectural signal.

\item \textbf{Run C1 (text-only) and C4 (multimodal) on N=108.}
Same H3/H4 evaluation set as Modification~2. Running the same two conditions
produces a directly comparable row in Table~\ref{tab:intrinsic}. If API budget
allows, run all four conditions (C1--C4) to produce a full row.

\item \textbf{Verify the output parser on 5--10 samples before the full run.}
Qwen3.5-397B-A17B output format may differ from Qwen3-VL: check that
Verdict, Error Category, and Feedback fields parse correctly with thinking
mode off. Normalize verdict strings to lowercase (e.g.,
\texttt{"Partially Correct"} $\to$ \texttt{"partially correct"}).

\item \textbf{Run the same three evaluation scripts as Modification~2.}
\begin{itemize}
  \item Verdict accuracy: predicted vs.\ ground-truth verdict labels (N=108)
  \item Intrinsic metrics: PC Recall, FSR, VGS from per-example raw outputs
  \item Feedback quality: LLM-as-Judge using \texttt{analyze\_feedback\_by\_error\_type.py},
    \texttt{model=claude-sonnet-4-6}, \texttt{temperature=0}, N=108
\end{itemize}

\item \textbf{Interpret against the three-way decision matrix:}

\begin{center}
\small
\begin{tabular}{lll}
\toprule
Mod~2 (9B) C4 vs C1 & Mod~3 (397B) C4 vs C1 & Implication for Report 4 \\
\midrule
Resolves gap & --- & Early fusion at 9B is sufficient; fine-tune Qwen3.5-9B \\
Does not resolve & Resolves gap & Scale + early fusion required; target 27B dense or 397B MoE \\
Does not resolve & Does not resolve & Task-inherent; use routing + Verdict-Classification Head \\
\bottomrule
\end{tabular}
\end{center}

\end{enumerate}

\textbf{Results.} \textit{[To be filled after running N=108 on C1 and C4.
Update Table~\ref{tab:intrinsic} with PC Recall, FSR, VGS, Verdict, and Fdbk
Match rows for Qwen3.5-397B-A17B C1 and C4.]}

% FIGURE PLACEHOLDER (Mod 3 combined): Seven-bar grouped chart comparing
% Qwen3-VL-8B C1, Qwen3-VL-8B C4, Routing, Qwen3.5-9B C1, Qwen3.5-9B C4,
% Qwen3.5-397B C1, Qwen3.5-397B C4.
% Three clusters: PC Recall, FSR, Fdbk Match. This is the primary Report 3 figure.

\clearpage

\subsection{Qualitative Analysis and Examples}
\label{sec:qual}

\noindent\textit{\textbf{NOTE:} This subsection is a WIP and will be refined.}
\medskip

\points{2} Table~\ref{tab:qual} presents five representative failure examples from
per-example model outputs (8B, text-only condition, H3 evaluation set). Two failure
modes account for the majority of errors: \textbf{Verdict Inflation}---the model
upgrades Partially Correct to Correct, producing no feedback---and
\textbf{Hallucinated Correction}---the model downgrades Partially Correct to
Incorrect with confident but wrong feedback. Both modes share a root cause: without
figure verification, the model falls back on plausibility assessment rather than
evidence-based judgment.

\begin{table*}[!ht]
\centering
\scriptsize
\setlength{\tabcolsep}{3pt}
\begin{tabular}{@{}p{0.03\linewidth}p{0.09\linewidth}p{0.17\linewidth}p{0.16\linewidth}p{0.17\linewidth}p{0.18\linewidth}p{0.13\linewidth}@{}}
\toprule
\# & Error Type & Question & Student Answer & Ground Truth & Model Output & Failure Mode \\
\midrule
1 &
Factual &
Which model performs best on Ubuntu for text lengths 60--90 words? &
``The LSTM model performs best for that length range.'' &
\textbf{Partially Correct.} Right range, misread top model from table. Feedback should correct the misread value and name the actual winner. &
\textbf{Incorrect.} ``The Transformer model performs best, not LSTM.'' &
\textbf{Hallucinated Correction.} Verdict wrong; feedback invents a specific correction that may itself be incorrect. Student given false confidence in wrong information. \\
\midrule
2 &
Factual &
Which of LSTw/oAR or LST-Skip performs better? &
``LSTw/oAR; its red forecast line stays closer to the true data than LST-Skip.'' &
\textbf{Partially Correct.} Correct comparison approach; misread which panel shows the closer match. &
\textbf{Correct. N/A.} &
\textbf{Verdict Inflation.} False positive; feedback suppressed. Student's panel misread goes uncorrected. \\
\midrule
3 &
Conceptual &
Which Ortiz et al.\ system had highest BLEU/METEOR, and what does this mean? &
``SMT system, 43.7 BLEU and 35.6 METEOR. This means SMT outperforms all systems.'' &
\textbf{Partially Correct.} Numbers correct, system identified correctly. Missed: interpretive claim is unsupported without cross-system comparison. &
\textbf{Correct. N/A.} &
\textbf{Verdict Inflation.} Numbers verified via pre-trained knowledge; interpretive gap not detected. Incomplete reasoning accepted as complete. \\
\midrule
4 &
Omission &
What is the role of temperature $t$ in CCA decoding? &
``Temperature $t$ controls the stopping condition; as $t$ decreases, the search ends.'' &
\textbf{Partially Correct.} Correctly notes $t$ decreases; omits its role shaping acceptance probabilities $\alpha_0, \alpha_1$ in the main loop. &
\textbf{Incorrect.} ``$t$ does not control stopping. It modulates a probability distribution over outputs.'' &
\textbf{Hallucinated Correction.} Partially-right claim over-penalized; actual omission (acceptance probabilities) not identified. Model contradicts a claim the student got right. \\
\midrule
5 &
Omission &
Which agent values the entire share $Z_j$? &
``Agent $j$ values $Z_j$, since the solid blue boxes cover all of $Z_j$.'' &
\textbf{Partially Correct.} Correct answer; omits that \emph{other} agents do not value the entire share---the contrastive point the question implies. &
\textbf{Correct. N/A.} &
\textbf{Verdict Inflation.} Plausible claim accepted without figure verification. Contrastive omission not detected. Most common mode for omission errors: 0/15 exact feedback match at C1. \\
\bottomrule
\end{tabular}
\caption{\points{2} Qualitative failure analysis (Qwen3-VL-8B, text-only condition, H3 evaluation
set). Two failure modes: \textbf{Verdict Inflation} (examples 2, 3, 5)---model predicts Correct
with N/A feedback, silencing the coaching pipeline; \textbf{Hallucinated Correction}
(examples 1, 4)---model predicts Incorrect with confident but wrong feedback. Root cause: without
the figure, the model performs plausibility checking rather than evidence-based
verification. Visual context (C4) reduces both modes by grounding reasoning in figure
content, explaining why feedback quality improves even as PC Recall drops.}
\label{tab:qual}
\end{table*}

\clearpage

\section{\points{3} Insights}

\noindent\textit{\textbf{NOTE:} This section is a WIP and will be refined.}
\medskip

\textbf{The verdict--feedback dissociation is a task decomposition signal, not a
bug to fix.} Our intrinsic metric analysis sharpens what appeared in Report 2 as a
surprising empirical finding into a precise, mechanistically-grounded claim. PC
Recall drops from 44.3\% to 19.7\% as visual input is added, while VGS rises from
35.5\% to 49.4\%. These measure different sub-skills that are genuinely in tension
at 8B scale: visual information improves value extraction into text (VGS) but
introduces noise into the three-class decision boundary, particularly for the
partially correct middle class that requires the most nuanced judgment. The
condition routing modification confirms this: using C1 for verdict and C4 for
feedback simultaneously recovers C1's PC Recall (44.3\%) and C4's feedback quality
(6.0\% match) at no additional inference cost. This reframes the main model
proposal for Report 4: rather than searching for a single configuration or fine-tuning
objective that improves both metrics jointly, Study Coach should be designed as an
explicit two-stage pipeline with task routing---a lightweight verdict classifier
(the Verdict-Classification Head proposed in Report 2, Section 3.4) decoupled from the
feedback generator. The head would be trained directly on the three-class boundary
using C1 or C2 inputs, while the feedback generator continues to use C4.

\textbf{The two architectural ablations together form a decision tree for Report 4.}
The gap between VGS (49.4\% at C4) and PC Recall (19.7\% at C4) has two competing
explanations. The first is architectural: Qwen3-VL's bolt-on SigLIP-2 encoder feeds visual
tokens into a Qwen3 backbone via a staged pipeline (merger-only Stage~0, then
full training), creating a representational boundary that may affect
classification more than generation. The second
is scale-dependent: the capacity required to use visual input reliably for three-class
verdict decisions simply exceeds what 8B parameters can provide, regardless of
architecture. Modifications~2 and 3 test these explanations at two operating points:
9B early fusion (isolates architecture at matched scale) and 397B early fusion (tests
whether the combination of architecture change and scale increase resolves what neither
alone could). The three outcomes map directly to Report 4 strategy: if Qwen3.5-9B resolves
the gap, fine-tune that backbone on SPIQA+; if only 397B resolves it, use the 27B
dense variant as a practical fine-tuning target (full 397B fine-tuning is cost-prohibitive
but the result identifies the correct architecture family); if neither resolves it, the
task requires the Verdict-Classification Head (Report 2, Section~3.4) as an explicit
supervisory signal for the verdict decision, with condition routing as the production
deployment strategy regardless. The condition routing result from Modification~1 is
robust to all three outcomes: it already captures C1 verdict precision and C4 feedback
quality simultaneously, and will remain the best zero-cost inference strategy until
fine-tuning is complete.

\subsection{Team Member Contributions to this Report}

\paragraph{Dipti Aswath} Defined the three intrinsic metrics (PC Recall, FSR, VGS)
and computed all values from per-example raw outputs. Designed and evaluated
Modification 1 (Condition Routing). Performed qualitative failure analysis and
selected representative examples. Wrote Sections 1 (Intrinsic Metrics),
Modification 1, Qualitative Analysis, Section 2 (Insights), and this contributions
section.

\paragraph{Buddhika Makehelwala} \textit{[To be filled in by Buddhika.]}

\paragraph{Jiashan Cui} \textit{[To be filled in by Jiashan.]}

\paragraph{Khubi Shah} \textit{[To be filled in by Khubi.]}

\bibliography{references}
\bibliographystyle{iclr2024_conference}

\appendix
\section{Appendix}

\noindent\textit{\textbf{NOTE:} This appendix is a WIP and will be refined.}
\medskip

\subsection{Intrinsic Metric Computation Details}

All intrinsic metrics computed from 8B per-example outputs (H3/H4 evaluation set,
N=108; 52 factual, 41 conceptual, 15 omission). Ground truth: 61 partially correct,
47 incorrect, 0 correct.

\textbf{PC Recall:} Count predictions of ``Partially Correct'' where ground truth =
``Partially Correct.'' 61 such cases in the dataset.

\textbf{FSR:} Count predictions where the raw output feedback field = ``N/A'' or
empty. These correspond exactly to Correct predictions by prompt design.

\textbf{VGS:} For factual error examples only, among non-suppressed feedbacks, count
those matching a regex for specific numerical values (\textit{digits followed by \%},
\textit{multi-digit integers}, or \textit{decimal numbers}). Denominators: C1=45,
C2=43, C3=36, C4=34 (remaining after suppression).

\textbf{Condition Routing:} Matched by (\texttt{paper\_id}, \texttt{question}) key
across text\_only and multimodal outputs. All 108 examples matched (100\% coverage).
Verdict and error category from text\_only; feedback from multimodal.

\subsection{Feedback Judgment Counts (8B, H3 Evaluation Set, N=108)}

\begin{table}[h]
\centering
\small
\begin{tabular}{@{}llrrrr@{}}
\toprule
Error Type & Condition & Match & Partial & Unmatched & Match\% \\
\midrule
Factual    & text\_only    &  7 &  9 & 36 & 13.5\% \\
Factual    & caption\_only & 13 &  7 & 32 & 25.0\% \\
Factual    & vision\_only  & 15 &  6 & 31 & 28.8\% \\
Factual    & multimodal    & 18 &  8 & 26 & 34.6\% \\
Conceptual & text\_only    & 11 & 16 & 14 & 26.8\% \\
Conceptual & caption\_only & 21 & 10 & 10 & 51.2\% \\
Conceptual & vision\_only  & 15 & 10 & 16 & 36.6\% \\
Conceptual & multimodal    & 24 &  4 & 13 & 58.5\% \\
Omission   & text\_only    &  0 &  7 &  8 &  0.0\% \\
Omission   & caption\_only &  0 &  7 &  8 &  0.0\% \\
Omission   & vision\_only  &  4 &  3 &  8 & 26.7\% \\
Omission   & multimodal    &  5 &  3 &  7 & 33.3\% \\
\bottomrule
\end{tabular}
\caption{LLM-as-Judge feedback evaluation counts (8B, N=108).}
\end{table}

\end{document}
